\section*{Acknowledgements}
In the course of this work, I profited from
many discussions with Filippo Palombi and Stefan Schaefer
of various questions related to
the slow topology-switching in current lattice QCD simulations.
I also wish to thank
Stefan Schaefer, Rainer Sommer and Francesco Virotta for
sharing some of their simulation results before publication.

\appendix

\numberwithin{equation}{section}

\section{SU(3) notation}

\subsection{Group generators}

The Lie algebra $\Lie$ of $\Group$
may be identified with the space of all anti-hermitian
traceless $3\times3$ matrices.
With respect to a
basis $T^a$, $a=1,\ldots,8$, of such matrices,
the elements $X\in\Lie$ are given by
\begin{equation}
  X=X^aT^a,
\end{equation}
where $(X^1,\ldots,X^8)\in\rz^8$
(repeated group indices are automatically summed over).

The generators $T^a$ are assumed to satisfy
the normalization condition
\begin{equation}
  \tr\bigl\{T^aT^b\bigr\}=-\frac{1}{2}\delta^{ab}.
\end{equation}
The structure of the Lie algebra is then encoded in the commutators
\begin{equation}
  [T^a,T^b]=f^{abc}T^c,
\end{equation}
while the completeness of the generators implies
\begin{align}
  \bigl\{T^a,T^b\bigr\}&=-\frac{1}{3}\delta^{ab}+id^{abc}T^c,
  \\[1.0ex]
  T^a_{\alpha\beta}T^a_{\gamma\delta}&=
  -\frac{1}{2}\left\{
  \delta_{\alpha\delta}\delta_{\beta\gamma}-\frac{1}{3}\delta_{\alpha\beta}
  \delta_{\gamma\delta}\right\}.
\end{align}
It follows from these equations that the structure constants
$f^{abc}$ and the tensor $d^{abc}$ are both real.
Moreover, $f^{abc}$
is totally anti-symmetric in the indices and $d^{abc}$ totally symmetric
and traceless.


\subsection{Adjoint representation}

The representation space of the adjoint representation of $\Lie$
is the Lie algebra itself,
i.e.\ the elements $X$ of $\Lie$ are represented by linear
transformations
\begin{align}
  &\Ad X:\,\Lie\mapsto\Lie,
  \\[0.8ex]
  &\Ad X\cdot Y=[X,Y]\quad\text{for all}\quad Y\in\Lie.
\end{align}
The action of $\Ad X$ on the group generators
is given by
\begin{equation}
  \Ad X\cdot T^b=T^a(\Ad X)^{ab},
\end{equation}
where
\begin{equation}
  (\Ad X)^{ab}=-f^{abc}X^c
\end{equation}
is a real antisymmetric $8\times 8$ matrix.


\subsection{Matrix norms}

The natural scalar product in $\Lie$ is
\begin{equation}
  (X,Y)=X^aY^a=-2\,\tr\bigl\{XY\bigr\}.
\end{equation}
In particular, $\|X\|=(X,X)^{1/2}$ is a possible definition of
the norm of $X\in\Lie$.

Another useful matrix norm derives from the square norm
\begin{equation}
  \|v\|_2=\bigl\{|v_1|^2+|v_2|^2+|v_3|^2\bigr\}^{1/2}
\end{equation}
of complex colour vectors $v$. If $A$ is any complex $3\times3$ matrix,
one defines
\begin{equation}
  \|A\|_2=\max_{\|v\|_2=1}\|Av\|_2.
\end{equation}
This norm satisfies
\begin{equation}
   \|A+B\|_2\leq\|A\|_2+\|B\|_2,
   \qquad
   \|AB\|_2\leq\|A\|_2\|B\|_2,
\end{equation}
for all matrices $A,B$. Moreover, if $A$ is hermitian or
antihermitian, $\|A\|_2$ is equal to the maximum of the absolute
values of its eigenvalues.

\section{Properties of the SU(3) exponential function}

\subsection{Lipschitz bound}

For any $X,Y\in\Lie$, the relation
\begin{equation}
  \|\rme^X-\rme^Y\|_2=
  \|1-\rme^{-X}\rme^Y\|_2
\end{equation}
follows from the fact that $\rme^X$ is unitary.
Using the identity
\begin{equation}
  1-\rme^{-X}\rme^Y=
  \int_0^1\rmd s\,\rme^{-sX}(X-Y)\rme^{sY}
\end{equation}
and the subadditivity (A.13) of the norm,
the Lipschitz bound
\begin{equation}
  \|\rme^X-\rme^Y\|_2\leq\|X-Y\|_2
\end{equation}
is then obtained.


\subsection{Differential of the exponential map}

Let $X$ be an element of the Lie algebra $\Lie$.
A linear mapping $J(X):\Lie\mapsto\Lie$ is then defined by
\begin{equation}
  J(X)\cdot Y=\rme^{-X}\left.{\rmd\over\rmd t}\rme^{X+tY}\right|_{t=0}
  \quad\text{for all $Y\in\Lie$.}
\end{equation}
$J(X)$ is referred to as the differential of the exponential map.
Scaling the expo\-nents by a parameter $s$, as above,
one obtains the representation
\begin{equation}
  J(X)\cdot Y=
  \int_0^1\rmd s\,\rme^{-sX}Y\rme^{sX}
\end{equation}
and thus the expansion
\begin{equation}
  J(X)=1+\sum_{k=1}^{\infty}{(-1)^k\over(k+1)!}(\Ad X)^k
\end{equation}
which is absolutely convergent for any $X\in\Lie$.

The action of $J(X)$ on the group generators $T^a$
is given by
\begin{equation}
  J(X)\cdot T^b=T^aJ(X)^{ab},
\end{equation}
where $J(X)^{ab}$ is a real $8\times8$ matrix.
Note that
\begin{equation}
  J(X)^T=J(-X)=\rme^{\Ad X}J(X),
  \qquad
  J(X)\Ad X=1-\rme^{-\Ad X}.
\end{equation}
Moreover, eq.~(B.5) implies
\begin{equation}
  \|J(X)\cdot Y\|_2\leq\|Y\|_2
\end{equation}
for all $X,Y\in\Lie$.

\section{Jacobian of the Euler step}

The Jacobian matrix (5.11) of the Euler step
is equal to unity except for some non-zero elements
along the row $(y,\nu)=(x,\mu)$. Its determinant
therefore coincides with the determinant of the $(x,\mu;x,\mu)$-element
\begin{equation}
  A^{ac}=\left\{\left(\rme^{\Ad X}\right)^{ac}
  +\eps J(-X)^{ab}[Z_{*}(U)](x,\mu;x,\mu)^{bc}
  \right\}_{X=\eps[Z(U)](x,\mu)}
\end{equation}
of the matrix.

The derivative $[Z_{*}(U)](x,\mu;x,\mu)^{bc}$
can be worked out explicitly in terms of
the plaquette sum
\begin{multline}
  M=\sum_{\nu\neq\mu}\bigl\{
  U(x,\mu)U(x+\hat{\mu},\nu)U(x+\hat{\nu},\mu)^{-1}U(x,\nu)^{-1}+\\
  U(x,\mu)U(x+\hat{\mu}-\hat{\nu},\nu)^{-1}U(x-\hat{\nu},\mu)^{-1}
  U(x-\hat{\nu},\nu)\bigr\}.
\end{multline}
A few lines of algebra then lead to the expression
\begin{equation}
  A^{ac}=B^{ac}+\frac{1}{2}\eps C^{ab}\Bigl\{
  id^{bcd}\tr\bigl\{T^d(M+M^{\dagger})\bigr\}
  -\frac{1}{3}\delta^{bc}\tr\bigl\{M+M^{\dagger}\bigr\}\Bigl\}
\end{equation}
where
\begin{align}
  B&=\frac{1}{2}\bigl(\rme^{\Ad X}+1\bigr),
  \\[0.6ex]
  C&=J(-X),\qquad X=-\eps\Proj\bigl\{M\bigr\}.
\end{align}
In particular,
\begin{equation}
  \det A=1-\frac{4}{3}\eps\,\tr\bigl\{M+M^{\dagger}\bigr\}+\rmO(\eps^2),
\end{equation}
as expected from eq.~(3.9).

\section{Inversion of the Euler step}

\subsection{Basic norm bounds}

For any complex $3\times3$ matrix $M$,
the inequality
\begin{equation}
  \|\Proj\bigl\{M\bigr\}\|_2\leq\frac{4}{3}\|M\|_2
\end{equation}
can be established in a few lines. One first observes that
\begin{equation}
  \frac{1}{2}\bigl(M-M^{\dagger}\bigr)=A D A^{-1},
  \qquad A\in\Group,
\end{equation}
where $D$ is a diagonal matrix with diagonal elements
$\lambda_1,\lambda_2,\lambda_3$. Setting
\begin{equation}
  \bar{\lambda}=\frac{1}{3}\sum_{k=1}^3\lambda_k,
\end{equation}
the estimates
\begin{multline}
  \|\Proj\bigl\{M\bigr\}\|_2=\max_k
  \left|\lambda_k-\bar{\lambda}\right|
  \\
  \leq\frac{4}{3}\max_k\left|\lambda_k\right|=
  \frac{2}{3}\|M-M^{\dagger}\|_2\leq\frac{4}{3}\|M\|_2
\end{multline}
then show that the inequality (D.1) holds for all matrices $M$.

An immediate consequence of the bound (D.1) and the Lipschitz bound (B.3)
is that
\begin{equation}
  \bigl\|\Proj\bigl\{A\bigl(\rme^X-\rme^Y\bigr)B\bigr\}\bigr\|_2\leq
  \frac{4}{3}\|X-Y\|_2
\end{equation}
for all $A,B\in\Group$ and $X,Y\in\Lie$.


\subsection{Solution of eq.~(5.1)}

For a given link $(x,\mu)$ and any fixed $U'(x,\mu)\in\Group$,
the function
\begin{equation}
  f(X)=\bigl\{[Z(U)](x,\mu)\bigr\}_{
  U(x,\mu)=\rme^{-\eps X}U'(x,\mu)}
\end{equation}
maps $X\in\Lie$ back to $\Lie$. Recalling eq.~(5.3), the inequality
(D.5) immediately implies that
\begin{equation}
  \|f(X)-f(Y)\|_2\leq k\|X-Y\|_2,
  \qquad k=8|\eps|.
\end{equation}
The function $f(X)$ is therefore a strict contraction if
the integration step size $\eps$ is in the range (5.5) (which is assumed
to be the case from now on).

It is not difficult to prove that
strict contractions in a complete metric space
have a unique fixed point
(see ref.~\cite{ReedSimonI}, Theorem V.18, for example).
In the present case, the fixed point
$X_{*}$ can be computed by noting that
the sequence $X_0=0,X_1,X_2,\ldots$ generated through
the recursion $X_{n+1}=f(X_n)$ satisfies
\begin{equation}
  \|X_n-X_{*}\|_2\leq k\|X_{n-1}-X_{*}\|_2
\end{equation}
and therefore rapidly converges to $X_{*}$.
The matrix
\begin{equation}
  U(x,\mu)=\rme^{-\eps X_{*}}U'(x,\mu)
\end{equation}
then provides a solution of eq.~(5.1).
Moreover, there is no other solution,
because the fixed point $X_{*}$ of $f(X)$ is unique.


\subsection{Positivity of the Jacobian}

In order to prove that
the determinant of the Jacobian matrix (C.1) is positive
at all $\eps$ in the range (5.5),
it suffices to show that the matrix has no zero mode, i.e.\ that
\begin{equation}
  A^{ac}Y^c=0
\end{equation}
implies $Y=0$.
To this end, eq.~(D.10) is first written in the form
\begin{equation}
  Y^a=-\eps J(X)^{ab}W^b,
\end{equation}
where
\begin{equation}
  X=\eps[Z(U)](x,\mu),\qquad
  W^b=[Z_{*}(U)](x,\mu;x,\mu)^{bc}Y^c.
\end{equation}
Recalling eq.~(5.10), the formula
\begin{equation}
  W=\lim_{t\to0}{1\over t}\left\{
  \left.[Z(U)](x,\mu)\right|_{U(x,\mu)\to\rme^{tY}U(x,\mu)}-[Z(U)](x,\mu)
  \right\}
\end{equation}
may then be derived from which one infers that
\begin{equation}
  \|W\|_2\leq 8\|Y\|_2.
\end{equation}
The inequality (D.5) has here been used again and also the fact that
the norm is a continuous map from $\Lie$ to $\rz$.
The combination of eq.~(D.11) and the bounds (B.9) and (D.14)
now implies $\|Y\|_2\leq k\|Y\|_2$ and thus $Y=0$.

\section{Differentiability of $\Sflow_t$}

The action $\Sflow_t$ solves an
elliptic partial differential equation with smooth coefficients
and is therefore guaranteed (by elliptic regularity)
to be a differentiable function of the
gauge field $U$.
In this appendix, the simultaneous
differentiability in the time $t$
is established by expanding $\Sflow_t$
in powers of $t-t_0$ around any fixed time $t_0$.
Using Sobolev norms,
the expansion can be shown to converge
if $t$ is sufficiently close to $t_0$.
The pointwise
uniform convergence of the series and its derivatives,
and therefore the differentiability of $\Sflow_t$,
then follows from
Sobolev's lemma.


\subsection{Sobolev spaces on the field manifold}

The definition of the Sobolev spaces on a compact manifold
is quite involved and will not be reviewed here. An introduction
to the subject and the proof of all statements made in this subsection
is given in the first three sections of ref.~\cite{Gilkey},
for example.

Let $\Cinfty$ be the space of differentiable functions on the field
manifold and $H_k$ the associated Sobolev space of order $k\in\gz$. The latter
is the completion of $\Cinfty$ with respect to a certain norm $\|\cdot\|_k$.
A characteristic feature of these norms is that
the bounds
\begin{equation}
  \|\mathfrak{D}\phi\|_k\leq c_k\|\phi\|_{k+p}
\end{equation}
hold for all $\phi\in\Cinfty$ and
any differential operator $\mathfrak{D}$ of order $p$
with smooth coefficients,
where the constants $c_k$ depend on $\mathfrak{D}$ but not on $\phi$.
Such differential operators thus extend to
bounded linear operators from $H_{k+p}$ to $H_k$.
Moreover,
$\|\phi\|_k\leq\|\phi\|_l$ if $k<l$
and therefore
$H_l\subset H_k$.

A fairly concrete description of the
Sobolev space $H_k$ can be given when $k$ is even
and non-negative. For any $t\in\rz$ and $j=0,1,2,\ldots$,
another norm $\|\phi\|_{t,j}$ of $\phi\in\Cinfty$
may be defined through
\begin{equation}
  \|\phi\|_{t,j}=\|\phi\|+\|\Lop^j\phi\|,
\end{equation}
where $\|\cdot\|$ is the norm associated to the scalar product (4.7).
The fact that $\Lop$ is a second-order elliptic differential operator
then implies
\begin{equation}
  \|\phi\|_{2j}\leq a_{t,j}\|\phi\|_{t,j},
  \qquad
  \|\phi\|_{t,j}\leq b_{t,j}\|\phi\|_{2j}
\end{equation}
for some constants $a_{t,j},b_{t,j}$.
In particular,
$H_{2j}$ is the completion of $\Cinfty$ with respect to the
norm $\|\cdot\|_{t,j}$.


\subsection{Properties of the inverse of $\Lop$}

Let $\Pop$ be the orthogonal projector to the zero mode of $\Lop$
in the Hilbert space with scalar product (4.7).
Its action on any function $\phi\in\Cinfty$ is given by
\begin{equation}
  \Pop\phi=(1,\phi)/(1,1).
\end{equation}
Note that $\Pop$ projects $\phi$ to a constant, but the constant
depends on the definition of the scalar product $(\cdot\,,\cdot)$
and therefore on $t$.
The action of the inverse $\Gop$ of $\Lop$ on $\phi$
is then determined by the equations
\begin{equation}
  \Lop\Gop\phi=(1-\Pop)\phi,
  \qquad
  \Pop\Gop\phi=0.
\end{equation}
As already mentioned in subsect.~4.2, the ellipticity of $\Lop$
and the absence of further zero modes
imply that these equations have, for any fixed $t$,
a unique differentiable solution $\Gop\phi$.
Evidently, $\Sflow_t=\Gop S$.

Since $\Lop$ has a spectral gap in the subspace orthogonal to
the zero mode, $\Gop$ is a bounded operator with respect
to the norm $\|\cdot\|$.
As a consequence, there exists a constant $g_t$ such that
\begin{equation}
  \|\Gop\phi\|_{t,j}\leq g_t\|\phi\|_{t,j-1}
\end{equation}
for all $j=1,2,\ldots$ and all $\phi\in\Cinfty$. Moreover,
recalling the inequalities (E.3), one concludes that
\begin{equation}
  \|\Gop\phi\|_{2j}\leq g_{t,j}\|\phi\|_{2j-2}
\end{equation}
for some other constants $g_{t,j}$.


\subsection{Expansion of $\Sflow_t$}

For any fixed time $t_0$, the operator $\Lop$ may be
decomposed according to
\begin{equation}
  \Lop=\Lopz+\left(t-t_0\right)\Lopp,
  \qquad
  \Lopp=\sum_{x,\mu}
  \left(\du^a_{x,\mu}S\right)\du^a_{x,\mu}.
\end{equation}
Each term in the power series
\begin{equation}
  \psi_t=\sum_{n=0}^{\infty}\left(t_0-t\right)^n
  \left(\Gopz\Lopp\right)^n\Gopz S
\end{equation}
is then a well-defined differentiable function of both $t$ and $U$.
Moreover, since
\begin{multline}
  \|\Gopz\Lopp\phi\|_{2j}\leq g_{t_0,j}\|\Lopp\phi\|_{2j-2}\leq
  d_jg_{t_0,j}\|\phi\|_{2j-1}\\
  \leq d_jg_{t_0,j}\|\phi\|_{2j},
\end{multline}
all $j=1,2,\ldots$ and some constants $d_j$, it is clear that the
series and all its derivatives with respect to $t$
converge uniformly in the norm $\|\cdot\|_{2j}$
if $t$ is sufficiently close to $t_0$.
Sobolev's lemma (statement (c) in lemma 1.3.5 of
ref.~\cite{Gilkey}) then implies
that the series converges pointwise and uniformly, together
with all its derivatives in $t$ and its derivatives in $U$ up
to some order proportional to $j$.

If $t$ is in a neighbourhood of $t_0$,
where the convergence of the
series and its derivatives up to order $l\geq2$
is guaranteed, the action of the
operator $\Lop$ and the summation may be interchanged
and one finds that
\begin{equation}
  \Lop\psi_t=S-\sum_{n=0}^{\infty}\left(t_0-t\right)^n
  \Popz\left(\Lopp\Gopz\right)^nS
  =\left(1-\Pop\right)S,
\end{equation}
the second equality being implied by the identities $\Pop\Lop\psi_t=0$
and $\Pop\Popz=\Popz$.
As a consequence,
\begin{equation}
  \left(1-\Pop\right)\psi_t=\Gop S=\Sflow_t,
\end{equation}
which shows that $\Sflow_t$ is,
in the specified neighbourhood of $t_0$,
$l$ times continuously differentiable with respect to $t$ and $U$.
Since $t_0$ and $l$ can be chosen arbitrarily, the
simultaneous differentiability
of $\Sflow_t$ is thus guaranteed at all times $t$ and to all orders.

\begin{thebibliography}{13}
\setlength{\itemsep}{0.2ex}

\bibitem{NicolaiMap}
H. Nicolai,
Phys. Lett. 89B (1980) 341;
Nucl. Phys. B176 (1980) 419

\bibitem{MoserZehnder}
J. Moser, E. J. Zehnder,
{\itshape Notes on Dynamical Systems},
Courant Lecture Notes, vol.~12 (American Mathematical Society, 2005)

\bibitem{DelDebbioTauQ}
L. Del Debbio, H. Panagopoulos, E. Vicari,
JHEP 0208 (2002) 044

\bibitem{StefanConference}
S. Schaefer, R. Sommer, F. Virotta,
{\itshape Investigating the critical slowing down of QCD simulations},
talk given at the XXVII International Symposium on Lattice Field Theory,
Beijing, China (July 2009), arXiv:0910.1465 [hep-lat],
to appear in the proceedings

\bibitem{HMC}
S. Duane, A. D. Kennedy, B. J. Pendleton, D. Roweth,
Phys. Lett. B195 (1987) 216

\bibitem{DuaneEtAl}
S. Duane, R. Kenway, B. J. Pendleton, D. Roweth,
Phys. Lett. B176 (1986) 143;
S. Duane, B. J. Pendleton,
Phys. Lett. B206 (1988) 101

\bibitem{ArnoldI}
V. I. Arnold,
{\itshape Ordinary Differential Equations}, 3rd ed.\ (Springer-Verlag, Berlin, 2008)

\bibitem{Gilkey}
P. B. Gilkey,
{\itshape Invariance Theory, the Heat Equation and the Atiyah-Singer
Index Theorem}, 2nd ed.\ (CRC Press, Boca Raton, 1995)

\bibitem{Wilson}
K. G. Wilson,
Phys. Rev. D10 (1974) 2445

\bibitem{HairerEtAl}
E. Hairer, C. Lubich, G. Wanner,
{\itshape Geometric Numerical Integration:
Structure-Preserving Algorithms for Ordinary Differential Equations},
2nd ed.\ (Springer, Berlin, 2006)

\bibitem{Stout}
C. Morningstar, M. Peardon,
Phys. Rev. D69 (2004) 054501

\bibitem{DDHMC}
M. L\"uscher,
Comput. Phys. Commun. 165 (2005) 199

\bibitem{ReedSimonI}
M. Reed, B. Simon,
{\itshape Methods of Modern Mathematical Physics}, vol.~I
(Academic Press, New York, 1972)

\end{thebibliography}
